% !TEX root =/main.tex

\chapter{Process scheduling}

\section{Introduction to process scheduling}
There's two different kinds of processes; CPU-bound and I/O bound.
\subsection{CPU-bound}
A cpu bound thread is a thread that is mainly using computing and not I/O devices, the main throttle is the speed of the processor not the access time of devices.

\subsection{I/O bound}
I/O bound threads are threads which perform a lot of I/O actions and therefore have to wait a lot inbetween when it can actually compute.

\section{Batch / Interactive / Rts scheduling}

\subsection{Batch}
''business-world” applications, data analysis. Appropriate for non-preemptive.

The most important metric is throughput.

\subsubsection{Comparison of different algorithms}
\begin{tabular}{|c|c|c|}
	\hline
	Algorithm                                 & +                               & -                                                   \\\hline
	First-Come/First-Serve (non-preemptive)   & Easy, fair                      & Possibly inefficient (esp. for I/O bound processes) \\\hline
	Shortest Job First (non-preemptive)       & Optimal for turnaround          & Starvation + need to know runtime                   \\\hline
	Shortest Remaining Time Next (preemptive) & New short jobs get good service & Starvation + need to know runtime                   \\\hline
\end{tabular}


\subsection{Interactive}
It is not known in advance what the system is  going to do. Many users that need responsiveness, requiring preemptive scheduling.

Most important metric is response time.

\subsubsection{Comparison of different algorithms}
\begin{tabular}{|c|c|c|}
	\hline
	Algorithm           & +                                  & -                                                \\\hline
	Round-robin         & Easy to implement                  & Easy to get too long or short times per process. \\\hline
	Priority scheduling & Can maximize the throughput easily & Harder to implement                              \\\hline
	Lottery scheduling  & High control over resource use     &                                                  \\
\end{tabular}

\subsection{Real-time}
We will not discuss this during the course, but if these do not run instantly, there is no point in running it. Think car control systems.

for short-lived, short-cycle processes with hard/soft deadlines.

Most important metric is meeting deadlines.

\section{Processes vs. Threads scheduling}

\section{Multiprocessor hardware}

\section{Multiprocessor scheduling}